\chapter{Proof of the non-collapsing}
\label{chap:proof-noncollapsing}
\label{sectnoncoll}

\section{The non-collapsing statement}

\begin{proposition}[Non-collapsing result assuming the strong canonical neighborhood assumption]
\label{prop:kappa-limit-noncollapsing}
\uses{thm:main-existence-ricci-flow-surgery}
\label{KAPPALIMIT}
For surgery parameter sequences through index $i$, there are $0<\kappa\leq\kappa_i$ and, for every $0<r_{i+1}\leq r_i$, a number $0<\delta(r_{i+1})\leq\delta_i$ with the following property. If $\overline\delta(t)\leq\delta_j$ on $[T_j,T_{j+1})$ and $\overline\delta(t)\leq\delta(r_{i+1})$ on $[T_{i-1},T_{i+1})$, then every Ricci flow with surgery on $[0,T)$, $T\in(T_i,T_{i+1}]$, satisfying Theorem~\ref{MAIN}, Assumptions~(1)--(7), and the strong $(C,\epsilon)$-canonical-neighborhood assumption with parameter $r_{i+1}$ is $\kappa$-noncollapsed on scales $\leq\epsilon$.
\end{proposition}
\begin{proof}
\sketch
The small-scale case follows from Proposition~\ref{prop:noncollapsing-small-scale-easy-case}. For larger scales, the cap-exclusion and compactness arguments below produce minimizing $\mathcal L$-geodesics to a controlled earlier ball; reduced-volume monotonicity then gives the uniform non-collapsing constant. The final controlled-ball construction is recorded after Claim~\ref{tau0cl}.
\end{proof}

\begin{remark}[Dependence of the non-collapsing constants and the two-outcome alternative]
\label{rem:kappa-delta-dependence}
\uses{prop:kappa-limit-noncollapsing}

The constants may also depend on the fixed data $t_0,\epsilon,C$, and $i+1$. At a point satisfying the non-collapsing hypotheses, either the relevant ball has the asserted volume lower bound or its time-slice component has positive sectional curvature.
\end{remark}

\section{Small scales}
\label{rsmall}

\begin{proposition}[Non-collapsing at scales $r\leq r_{i+1}$ follows directly from canonical neighborhoods]
\label{prop:noncollapsing-small-scale-easy-case}
\uses{def:strong-canonical-neighborhood-assumption}

If ${\bf t}(x)=t$ and $R(x)=r^{-2}\geq r_{i+1}^{-2}$, then there is $\kappa>0$, depending only on $C$, such that either $\Vol B(x,t,r)\geq\kappa r^3$ or $x$ belongs to a positively curved component of $M_t$.
\end{proposition}
\begin{proof}
\sketch The strong $(C,\epsilon)$-canonical neighborhood is a neck, an $\epsilon$-round component, a $C$-component, or a cap. The first has a universal volume ratio and the middle two give the positive-curvature alternative. In the cap case, choose $r'$ with $\sup_{B(x,t,r')}|\Rm|=(r')^{-2}$. Cap control gives $\Vol B(x,t,r')\geq C^{-1}(r')^3$ and $r'/r\geq C^{-1/2}$, hence $\Vol B(x,t,r)\geq C^{-5/2}r^3$.
\end{proof}

\section{Minimizing $\mathcal L$-geodesics}

\begin{proposition}[Existence and compactness of minimizing $\mathcal L$-geodesics avoiding the surgeries]
\label{prop:minimizing-L-geodesics-exist}
\uses{prop:kappa-limit-noncollapsing}
\label{delta0ri+1}
For every $0<r_{i+1}\leq r_i$ there is $\delta=\delta(r_{i+1})>0$ such that, if $\overline\delta\leq\delta$ on $[T_{i-1},T_{i+1}]$, the following holds. Let $x$ have time $T\in[T_i,T_{i+1})$, and suppose $P(x,r,T,-r^2)$ exists with $r\geq r_{i+1}$ and $|\Rm|\leq r^{-2}$. There is an open smooth subset $U\subset{\bf t}^{-1}[T_{i-1},T)$ such that: every $y\in U$ is joined to $x$ by a minimizing $\mathcal L$-geodesic; $U_t\ne\varnothing$ for $t\in[T_{i-1},T)$; $\min_{U_t}\mathcal L\leq3\sqrt{T-t}$; and the minimizing locus with bounded $\mathcal L$-length is compact over compact time intervals.
\end{proposition}
\begin{proof}
\sketch
\uses{cor:minimizing-geodesic-to-earlier-slice,prop:L-geodesics-exist-and-limit}
The cap and disappearing-region exclusion lemmas below confine short paths to compact smooth sets. Their compactness and the reduced-length estimate establish the four stated properties.
\end{proof}

\section{Evolution of surgery caps}

\begin{proposition}[Evolution of a surgery cap stays close to the standard flow until removed]
\label{prop:surgery-cap-evolution-alternative}
\uses{def:surgery-cap, def:curvature-pinched-toward-positive-surgery}
\label{altern}
For $A<\infty$, $\delta''>0$, and $0<\theta<1$, there is $\delta''_0=\delta''_0(A,\theta,\delta'')$, also depending on the fixed $r_{i+1},C,\epsilon$, with the following property. If surgery at $\bar t$ has scale $\bar h$, tip $p$, and $\overline\delta(\bar t)\leq\delta''_0$, put $T'=\min(T,\bar t+\theta\bar h^2)$. Either $B(p,\bar t,A\bar h)$ evolves by a compatible embedding through $T'$, and after shifting and scaling by $\bar h^{-2}$ is $\delta''$-close in $C^{[1/\delta'']}$ to the standard flow, or it evolves by such an embedding to some $\bar t_+\in(\bar t,T')$ and is contained in $D_t$ for $t<\bar t_+$ sufficiently close to $\bar t_+$.
\end{proposition}
\begin{proof}
  \uses{prop:standard-flow-asymptotic-cylinder, thm:shi-derivative-estimates}
\sketch

Assume failure and rescale a counterexample sequence with control parameters tending to zero. The cap-limit results below give standard-flow closeness up to each maximal embedding time. If the embedding stops early, Claim~\ref{claim:ball-disjoint-from-surgery-spheres} and connectedness force all flow lines in the ball to disappear at that surgery; otherwise the ball evolves through $T'$.
\end{proof}

\begin{claim}[Uniform derivative bounds for the rescaled surgery-cap metric]
\label{claim:surgery-cap-curvature-bound-uniform}
If a surgery cap is attached to a strong $\delta'$-neck of scale $h$, then for every $\ell<\infty$ the rescaled initial metric on the middle neck half together with the cap has a uniform bound on $|\nabla^\ell\Rm|$.
\end{claim}
\begin{proof}
\sketch
\uses{thm:shi-derivative-estimates}
Rescaling gives a neck flow with bounded curvature on a unit backward interval, so Shi estimates control the middle half. The fixed smooth cap gluing converges to the standard initial cap as $\delta'\to0$, preserving uniform bounds.
\end{proof}

\begin{corollary}[Smooth limit of rescaled surgery caps is the standard initial metric]
\label{cor:surgery-cap-smooth-limit}
For $\delta'_n\to0$, surgery caps attached to strong $\delta'_n$-necks of scales $h_n$, rescaled by $h_n^{-2}$ and based at their tips, subconverge smoothly to the standard initial surgery-cap metric.
\end{corollary}
\begin{proof}
\sketch Apply Claim~\ref{claim:surgery-cap-curvature-bound-uniform}; the cap-gluing construction identifies the limit with the standard initial metric.
\end{proof}

\begin{lemma}[Rescaled Ricci flows with surgery near a surgery cap converge to the standard flow]
\label{lem:rescaled-surgery-limit-standard-flow}
Let surgery scales $h_n$ and controls $\delta'_n\to0$ occur at $t_n$, with cap tips $p_n$, and suppose compatible embeddings of every fixed rescaled ball persist until times $\theta_n\to\theta<1$. After shifting $t_n$ to $0$ and scaling by $h_n^{-2}$, a subsequence converges smoothly on $[0,\theta)$ to the standard flow based at its tip.
\end{lemma}
\begin{proof}
\sketch
\uses{thm:shi-with-initial-derivatives}
The preceding corollary gives the standard initial limit. On each shorter interval, uniqueness identifies the limit with the standard flow; curvature control, positive pinching, and Shi estimates extend this by a uniform time. A diagonal subsequence reaches $[0,\theta)$.
\end{proof}

\begin{corollary}[Surgery-cap evolution is close to the standard flow up to time $\theta$]
\label{cor:surgery-cap-close-to-standard-up-to-theta}\label{goodtotheta}
Under the hypotheses of Lemma~\ref{lem:rescaled-surgery-limit-standard-flow}, for every $A<\infty$ and $\delta''>0$, the rescaled flow on $B(p_n,0,A)\times[0,\theta_n)$ is $\delta''$-close in $C^{[1/\delta'']}$ to the standard flow for all sufficiently large $n$.
\end{corollary}
\begin{proof}
\sketch
\uses{thm:shi-with-initial-derivatives}
Smooth convergence holds on every $[0,\theta-\eta']$. Uniform curvature and derivative bounds on the remaining short interval, followed by $\eta'\downarrow0$, give the stated closeness.
\end{proof}

\begin{claim}[Surgery spheres stay disjoint from the ball around the surgery cap]
\label{claim:ball-disjoint-from-surgery-spheres}
If the maximal cap embedding stops before $\min(\theta,(T'_n-t_n)/h_n^2)$, then for sufficiently small $\delta'_n$ its image is disjoint, near the stopping time, from the backward images of every surgery sphere at that time.
\end{claim}
\begin{proof}
\sketch
\uses{cor:surgery-cap-close-to-standard-up-to-theta,cor:embedded-sphere-isotopic-to-neck-fiber}
Closeness to the standard flow bounds curvature and diameter in units of $h_n$. An intersection would force the cap ball into a much longer later surgery neck. Its internal neck sphere would then be isotopic to the central sphere of that neck, although it bounds a $3$-ball, a contradiction.
\end{proof}

\begin{remark}[A surgery-cap neighborhood may later be entirely removed by a further surgery]
\label{rem:ball-removed-by-later-surgery}
\uses{prop:surgery-cap-evolution-alternative}

A cap ball in the preceding proposition can be removed later, for instance as part of a capped $\epsilon$-horn.
\end{remark}

\section{Length estimates near surgery}

\begin{definition}[The constant $\overline\delta_0$ controlling closeness to the standard metric]
\label{def:overline-delta-0}\label{defnoverlinedelta0}
\uses{prop:standard-solution-scalar-curvature-lower-bound}

Let $c>0$ be the constant of Proposition~\ref{Restim}. Fix $0<\overline\delta_0<1/4$ so that $C^{[1/\overline\delta_0]}$-closeness to the standard initial metric implies $|R_{g'}-R_{g_0}|<c/2$ and $|\Ric_{g'}-\Ric_{g_0}|<1/4$.
\end{definition}

\begin{proposition}[Paths of bounded energy near a surgery cap have long integral length]
\label{prop:long-L-length-near-surgery-cap}
\uses{prop:surgery-cap-evolution-alternative}
\label{bigell}
For every $\ell<\infty$ there are $A_0=A_0(\ell)$ and $0<\theta_0=\theta_0(\ell)<1$ such that, for $A\geq A_0$ and sufficiently small surgery control, the following holds. Let a surgery cap of scale $h$ and tip $p$ evolve on $P(p,\bar t,Ah,T'-\bar t)$, where $T'=\min(T,\bar t+\theta_0h^2)$. A backward-time curve in this closure from the final slice or its outer boundary to $B(p,\bar t,Ah/2)$ at a time $t'\in[\bar t,\bar t+h^2/2]$ satisfies
\[
\int_0^{T'-t'}\bigl(R(\gamma(\tau))+|X_\gamma(\tau)|^2\bigr)\,d\tau>\ell.
\]
\end{proposition}
\label{*}\label{**}
\begin{proof}
  \uses{prop:surgery-cap-evolution-alternative, prop:standard-solution-scalar-curvature-lower-bound, prop:standard-flow-asymptotic-cylinder}
\sketch

Rescale to $\bar t=0$ and $h=1$. Choose $\theta_0$ close enough to $1$ that integrating $R\geq c/(2(1-t))$ exceeds $\ell$ for a curve reaching the final time. For the boundary alternative, the cylindrical separation and metric comparison below yield energy at least $A^2/75$; take $A_0\geq10\sqrt\ell$.
\end{proof}

\begin{claim}[Standard solution neck structure far from the tip persists up to time $\theta_0$]
\label{claim:standard-solution-neck-at-infinity}
For some $A'_0<\infty$, every $A\geq A'_0$ has the property that the standard flow on $(B(p_0,0,A)\setminus B(p_0,0,A/2))\times[0,\theta_0]$ is close to an evolving cylinder. In particular, its time-$t$ distance from $B(p_0,0,A/2)$ to the complement of $B(p_0,0,A)$ exceeds $A/4$.
\end{claim}
\begin{proof}
\sketch
\uses{prop:standard-flow-asymptotic-cylinder}
This is the asymptotic-cylinder estimate for the standard flow, uniformly on the compact interval $[0,\theta_0]$.
\end{proof}

\begin{lemma}[Curves with small $\mathcal L_+$-length avoid neighborhoods of surgery caps]
\label{lem:small-L-plus-avoids-surgery-caps}
\uses{prop:long-L-length-near-surgery-cap}
\label{L_0}
Given $L_0<\infty$, there is $\overline\delta_1=\overline\delta_1(L_0,r_{i+1})>0$ such that, if $\overline\delta\leq\overline\delta_1$ on $[T_{i-1},T)$, every backward curve from $x$ with $\mathcal L_+<L_0$ remains in $P(x,T,r/2,-r^2)$ for
\[
\tau\leq\min\left\{\tau_0,\frac{r_{i+1}^4}{256L_0^2},\log(\sqrt[3]2)r_{i+1}^2\right\},
\]
and avoids every compatible surgery-cap neighborhood $B(p,\bar t,(50+A_0)\bar h)\times[\bar t,t']$ with $t'\leq\bar t+\bar h^2/2$.
\end{lemma}
\label{tau_1est}
\begin{proof}
\sketch
\uses{prop:surgery-cap-evolution-alternative,prop:long-L-length-near-surgery-cap}
Choose $\ell=L_0/\sqrt{\tau'}$, then take $A$ and $\overline\delta_1$ from Proposition~\ref{bigell}. The velocity estimate in Claim~\ref{claim:energy-lower-bound-exit-ball} and Cauchy--Schwarz exclude an early exit. Any later entry into a cap neighborhood contradicts Proposition~\ref{bigell} after multiplication by $\sqrt{\tau'}$.
\end{proof}

\begin{remark}[Recollection: the surgery cap radius in terms of the surgery scale]
\label{rem:surgery-cap-radius-recollection}
\uses{def:surgery-cap}

The surgery cap attached at scale $\bar h$ has radius $(A_0+4)\bar h$ in the rescaled standard initial metric.
\end{remark}

\begin{claim}[Lower bound on the velocity integral before exiting a fixed ball]
\label{claim:energy-lower-bound-exit-ball}
If $\gamma$ first leaves $P(x,T,r/2,-r^2)$ at time $\tau''$ in the preceding lemma, then $\int_0^{\tau''}|X_\gamma|\,d\tau>r/(2\sqrt2)$.
\end{claim}
\begin{proof}
\sketch Curvature control gives $e^{-2r^{-2}\tau''}G(T)\leq G(T-\tau'')\leq e^{2r^{-2}\tau''}G(T)$, and the choice of $\tau''$ makes the distortion less than $2$. The final point has $G(T)$-distance $r/2$ from $x$.
\end{proof}

\section{Disappearing regions and compactness}

\begin{definition}[The disappearing region $J$ near a surgery time]
\label{def:disappearing-region-J}
\uses{def:surgery-operation}
\label{J}
Let $\bar t$ be a surgery time with no earlier surgery in $(\bar t-\tau_1,\bar t)$. Flow backward the neck portions $J_0(\bar t)=\coprod_i s_{N_i}^{-1}(-25,0]$ adjacent to its surgery spheres, and write $J_0(t)$ for their images. If $D_t$ is the region disappearing at $\bar t$, set $J(t)=J_0(t)\cup D_t$ and $J(\bar t-\tau_1,\bar t)=\bigcup_{t\in(\bar t-\tau_1,\bar t)}J(t)$.
\end{definition}

\begin{lemma}[Paths of bounded kinetic energy avoid the disappearing region]
\label{lem:short-energy-avoids-disappearing-region}
\uses{def:disappearing-region-J}
\label{15.17}
Fix the flow and $x$ as in Proposition~\ref{delta0ri+1}. For $1<\ell<\infty$, let $\bar t$ be a surgery time of scale $\bar h$, and suppose $0<\tau_1\leq\ell^{-1}\bar h^2$ contains no earlier surgery. A path beginning outside all $B(p_i,\bar t,(50+A_0)\bar h)$ and satisfying $\int_0^{\tau_1}|X_\gamma|^2\leq\ell$ is disjoint from $J(\bar t-\tau_1,\bar t)$.
\end{lemma}
\begin{proof}
\sketch The initial point is separated from the relevant neck boundary by order $\bar h$. Small surgery control keeps the neck metrics uniformly comparable during $\tau_1$. Entering $J$ therefore gives $\int|X_\gamma|>10\bar h$, and Cauchy--Schwarz gives $\int|X_\gamma|^2>\ell$.
\end{proof}

\begin{definition}[Compact subset $Y(\ell)$ of smooth points avoiding all singular regions]
\label{def:compact-subset-Y-ell}
\uses{def:disappearing-region-J, prop:long-L-length-near-surgery-cap}

Fix $\ell<\infty$. At each surgery time $\bar t\in[T_{i-1},T]$, remove the evolving cap neighborhoods of radius $(A_0+10)h(\bar t)$ through their maximal standard-flow time and remove $\widetilde J(\bar t-\tau_1(\ell,\bar t),\bar t)$, where $\tau_1(\ell,\bar t)=\min(h(\bar t)^2/\ell,\bar t-\bar t')$ and $\bar t'$ is the preceding surgery time or $0$. Let $Y(\ell)$ be the complement of the union of these open singular neighborhoods in ${\bf t}^{-1}([T_{i-1},T])$. Then $Y(\ell)$ is compact and lies in the smooth locus.
\end{definition}

\begin{proposition}[Existence, semicontinuity, and compactness of limits of short $\mathcal L$-length paths]
\label{prop:L-geodesics-exist-and-limit}
\uses{lem:small-L-plus-avoids-surgery-caps}
\label{Lgeoexist}
Fix $0<L<\infty$ and set $L_0=L+4(T_{i+1})^{3/2}$. If $\overline\delta\leq\overline\delta_1(L_0,r_{i+1})$ on $[T_{i-1},T_{i+1}]$, then any sequence of backward paths from $x$, of common duration at most $T-T_{i-1}$ and $\mathcal L$-length at most $L$, has a uniformly convergent subsequence with smooth-locus limit $\gamma$ satisfying
\[
\mathcal L(\gamma)\leq\liminf_n\mathcal L(\gamma_n).
\]
For fixed endpoint, a minimizing sequence limits to a minimizing $\mathcal L$-geodesic. Moreover, every backward path of $\mathcal L$-length at most $L$ into ${\bf t}^{-1}([T_{i-1},T))$ lies in $Y(\ell)$ for some $\ell$ depending only on $L$.
\end{proposition}
\label{L'ineq}\label{gamliminf1}\label{gamliminf2}
\begin{proof}
\sketch
\uses{lem:small-L-plus-avoids-surgery-caps,prop:surgery-cap-evolution-alternative,prop:long-L-length-near-surgery-cap,lem:l-geodesic-existence}
Claims~\ref{claim:path-misses-P0-cap-neighborhood} and~\ref{claim:path-disjoint-from-J-near-surgery} place short paths in $Y(\ell)$. Weighted energy bounds give equicontinuity and weak $W^{1,2}$ compactness on its smooth pieces; lower semicontinuity proves the limit statement, and a minimizing sequence gives the fixed-endpoint conclusion.
\end{proof}

\begin{claim}[A short path from $x$ misses the extended surgery-cap neighborhood $P_0$]
\label{claim:path-misses-P0-cap-neighborhood}
Any backward path from $x$ with $\mathcal L$-length less than $L$ is disjoint from the extended cap neighborhood $P_0(\mathcal C)$ associated to every surgery cap $\mathcal C$.
\end{claim}
\begin{proof}
\sketch
\uses{lem:small-L-plus-avoids-surgery-caps}
The lower scalar-curvature bound $R\geq-6$ changes the $\mathcal L$ bound to $\mathcal L_+\leq L+4(T_{i+1})^{3/2}=L_0$. Lemma~\ref{L_0} excludes the corresponding cap neighborhood.
\end{proof}

\begin{claim}[A short path from $x$ is disjoint from $J$ near each surgery time]
\label{claim:path-disjoint-from-J-near-surgery}
For every surgery time $\bar t\in[T_{i-1},T]$, every backward path from $x$ with $\mathcal L$-length at most $L$ is disjoint from $J(\bar t-\tau_1(\bar t),\bar t)$.
\end{claim}
\begin{proof}
\sketch
\uses{lem:short-energy-avoids-disappearing-region}
Claim~\ref{claim:path-misses-P0-cap-neighborhood} supplies the initial separation from cap tips. Away from the initial short interval, the weighted $\mathcal L_+$ estimate gives the kinetic-energy hypothesis of Lemma~\ref{15.17}; the initial interval is already disjoint by curvature separation.
\end{proof}

\begin{corollary}[Existence of a minimizing $\mathcal L$-geodesic to an earlier time-slice]
\label{cor:minimizing-geodesic-to-earlier-slice}
\uses{lem:small-L-plus-avoids-surgery-caps, prop:L-geodesics-exist-and-limit}
\label{bardelta1}
Given $L<\infty$, let $\overline\delta_1=\overline\delta_1(L+4T_{i+1}^{3/2},r_{i+1})$. If $\overline\delta\leq\overline\delta_1$ on $[T_{i-1},T_{i+1}]$, then every $\mathcal L$-path of length at most $L$ from $x\in{\bf t}^{-1}([T_i,T_{i+1}))$ to $y\in M_{T_{i-1}}$ admits a minimizing smooth-locus $\mathcal L$-geodesic from $x$ to $y$.
\end{corollary}
\begin{proof}
  \uses{cor:minimizing-geodesic-to-earlier-slice, prop:L-geodesics-exist-and-limit}
\sketch Apply Proposition~\ref{Lgeoexist} to a minimizing sequence.
\end{proof}

\section{Proof of Proposition~\ref{delta0ri+1}}

\begin{claim}[Compactness of the minimal-$\mathcal L$-length locus $Z'$ in each time-slice]
\label{claim:Z-prime-compact}
Let $Z$ be the locus minimizing $\mathcal L_x$ in its $U$ time-slice and put $Z'=\{z\in Z:\mathcal L_x(z)\leq L/2\}$. For every compact $I\subset[T_{i-1},T)$, $Z'\cap{\bf t}^{-1}(I)$ is compact.
\end{claim}
\begin{proof}
\sketch Take a convergent sequence of times and minimizing geodesics. Proposition~\ref{Lgeoexist} provides a smooth limiting geodesic. Flowing or restricting it across the limiting time proves that the limit realizes the limiting minimum, excluding a strict lower value by local extension in either time direction.
\end{proof}

\begin{lemma}[The minimum reduced length along the surviving interval is at most $3/2$]
\label{lem:reduced-length-min-bound-3-2}\label{lvalue}
If $l_x^{\mathrm{min}}(\tau)$ denotes the minimum of $l_x$ on $U_{T-\tau}$, then $l_x^{\mathrm{min}}(\tau)\leq3/2$ throughout the interval on which all preceding reduced lengths are at most $L/2$.
\end{lemma}
\begin{proof}
\sketch
\uses{cor:minimizing-geodesic-to-earlier-slice,prop:minimizing-L-geodesics-exist,cor:generalized-flow-n-half-bound}
Apply Corollary~\ref{compactcoro} using the already established existence, nonemptiness, and compactness properties of $U$.
\end{proof}

\section{Completion of the non-collapsing proof}
\label{Lgamma1}

\begin{claim}[Existence of a point of curvature below $r_i^{-2}$ along the $\mathcal L$-geodesic]
\label{claim:existence-of-tau0-low-curvature-point}\label{tau0cl}
Let $\gamma$ be an $\mathcal L$-geodesic from $x$ to $M_{T_{i-1}}$ with $\mathcal L(\gamma)\leq3\sqrt{T-T_{i-1}}$. There is $\tau_0\in[\max(\epsilon^2,T-T_i),T-T_{i-1}-\epsilon^2]$ such that $R(\gamma(\tau_0))<r_i^{-2}$.
\end{claim}
\begin{proof}
\sketch If $R\geq r_i^{-2}$ on the displayed interval, Claim~\ref{claim:tau0-estimates} gives an $\mathcal L$-lower bound of $4\sqrt{T-T_{i-1}}$, contradicting $\mathcal L(\gamma)\leq3\sqrt{T-T_{i-1}}$.
\end{proof}

\begin{claim}[Auxiliary time-interval estimates for the low-curvature point]
\label{claim:tau0-estimates}
For $T'=\max(\epsilon^2,T-T_i)$ and $T''=T-T_{i-1}-\epsilon^2$,
\[
(T'')^{3/2}-(T')^{3/2}\geq\tfrac14(T-T_{i-1})^{3/2},\quad
4(T')^{3/2}\leq4(T-T_{i-1})^{3/2},
\]
and
\[
4\bigl((T-T_{i-1})^{3/2}-(T'')^{3/2}\bigr)
\leq\frac{2t_0}{25}\sqrt{T-T_{i-1}}.
\]
\end{claim}
\begin{proof}
\sketch The bounds $T''/(T-T_{i-1})\geq0.9$, $T'/(T-T_{i-1})\leq2/3$, and $\epsilon^2\leq t_0/50$ give the three inequalities.
\end{proof}
